\begin{textsample}{POS Dim 4 – 2019\_c5 – Score 1.00 – 2019\_c5/enade\_2019\_zoot\_807802c04e.txt}  \label{ex:f4_pos_199}
CAPÍTULO 5 DISTRIBUIÇÃO DOS CONCEITOS 5.1 PANORAMA NACIONAL DA DISTRIBUIÇÃO DOS CONCEITOS O Conceito Enade é um indicador de qualidade que avalia os cursos por intermédio dos desempenhos dos estudantes a partir dos resultados do Enade.
Os valores possíveis vão de 1 ( pior situação ) a 5 ( melhor situação ).
Os cursos com apenas um ou sem nenhum concluinte participante não obtêm o Conceito Enade, ficando Sem Conceito ( SC ).
Na Tabela 5.1, são apresentadas a quantidade e a distribuição de cursos de Zootecnia participantes do Enade/2019, por faixa de conceito e Grande Região.
Enfatiza -se, mais uma vez, que a diferença entre os cursos tabulados neste capítulo e no capítulo 2 corresponde aos cursos que não foram avaliados, ou seja, àqueles, em princípio, sem alunos concluintes que tivessem sido inscritos para a prova.
Observando -se os dados da Tabela 5.1, nota -se que, dos 97 cursos participantes, 32 ( 33,0 \% ) classificaram -se com conceito 3, o valor modal.
Na região Norte, os conceitos modais foram 1 ou 3 ( 30,8 \% para cada ), o conceito 2 foi a classe modal para região Nordeste ( 30,8 \% ) e Centro-Oeste ( 44,4 \% ).
Na região Sudeste, os conceitos modais foram 3 e 4 ( 41,7 \% para cada ), e, na região Sul a classe modal foi o conceito 4 ( 56,3 \% ).
O conceito 4 foi o segundo mais frequente em nível nacional ( 28,9 \%, correspondendo a 28 cursos ), e o conceito 2, o terceiro ( 22,7 \%, correspondendo a 22 cursos ).
Houve, ainda, quatro cursos ( 4,1 \% ) que receberam conceito 5, e dez cursos ( 10,3 \% ) que receberam conceito 1.
Dos 96 cursos de Zootecnia, um ficou Sem Conceito ( SC ).
A região Norte participou com 13 cursos ou 13,4 \% do total nacional.
Desses, quatro cursos receberam o conceito 1 e quatro receberam conceito 3 ( 30,8 \% para cada ), além de ser o conceito modal na região, como já comentado.
Dois cursos ( 15,4 \% ) receberam conceito 2, e a três cursos ( 23,1 \% ), conceito 4.
Nenhum curso ficou Sem Conceito ( SC ) e nenhum curso recebeu o conceito 5 nessa região.
A região Nordeste participou com 26 cursos ou 26,8 \% do total nacional.
Nessa região, oito cursos ( 30,8 \% em termos regionais ) obtiveram conceito 2, o conceito modal para a região, como já comentado.
O conceito 1 foi atribuído a cinco cursos ( 19,2 \% ), o conceito 3 a sete cursos ( 26,9 \% ), conceito 4 a quatro cursos ( 15,4 \% ) e conceito 5 a um curso ( 3,8 \% ).
Um curso ficou Sem Conceito ( SC ).
Dos 24 cursos participantes da região Sudeste ( 24,7 \% do total nacional ), os conceitos 3 e 4 foram os conceitos modais, como já anteriormente assinalado, tendo sido obtido por dez cursos cada ( 41,7 \%, para cada ).
O conceito 2 foi atribuído a dois cursos ( 8,3 \% ), e o conceito 5 também foi atribuído a dois cursos ( 8,3 \% ).
Nessa região nenhum curso ficou Sem Conceito ( SC ) ou recebeu conceito 1.
Os 16 cursos da região Sul corresponderam a 16,5 \% do total nacional.
Como já indicado, houve predominância do conceito 4, atribuído a 56,3 \% dos cursos da região ( conceito modal ), ou seja, atribuído a nove dos 16 cursos participantes na região Sul.
O conceito 2 foi atribuído a dois cursos ( 12,5 \% ) e o conceito 3 a cinco cursos ( 31,3 \% ).
Nenhum curso ficou Sem Conceito ( SC ) na região Sul, e nenhum curso recebeu conceito 1 e 5.
Já dos 18 cursos participantes na região Centro-Oeste ( 18,6 \% do total nacional ), oito ( 44,4 \% em termos regionais ) receberam conceito 2, o conceito modal, como já fora destacado.
Um curso ( 5,6 \% ) obteve conceito 1, outros seis ( 33,3 \% ), conceito 3.
Além disso, o conceito 4 foi atribuído a 2 cursos ( 11,1 \% ) e conceito 5 a um curso ( 5,6 \% ).
Nenhum curso ficou Sem Conceito ( SC ).
No Gráfico 5.1, é apresentada a distribuição cumulativa do Conceito Enade segundo Grande Região.
Em um gráfico de distribuição cumulativa, em cada valor de abscissa ( eixo x ), o valor apresentado no eixo das ordenadas ( eixo y ) é o valor acumulado da variável ( conceito Enade ).
Observando a linha vertical pontilhada, que ajuda a visualizar o percentual de conceitos acumulados até o valor 3, ou seja, conceitos 1, 2 e 3, pode -se notar que, por exemplo, a região Centro-Oeste ( linha laranja ) apresenta 83 \% dos cursos com conceito 3 ou abaixo deste valor.
A região Sul ( linha rosa ), por outro lado, apresenta 44 \% dos cursos na mesma situação : com conceito 3 ou abaixo.
Destaca -se que em algumas regiões o conceito 5 não foi alcançado por nenhum curso, como ocorre nas regiões Norte ( linha azul ) e Sul ( linha rosa ), nas quais todos os cursos têm conceito menor ou igual a 4, e as respectivas poligonais alcançam 100 \% no conceito 4.
Em linhas gerais, em um gráfico de distribuição cumulativa, poligonais mais à esquerda representam Grandes Regiões com pior distribuição, e poligonais mais à direita, Grandes Regiões com melhores conceitos.
A região Norte ( linha azul ) apresenta os cursos com conceitos mais baixos, ao passo que as regiões Sudeste ( linha vermelha ) e Sul ( linha rosa ) apresentam as distribuições com valores maiores.
Cursos em IES situadas nas regiões Centro-Oeste ( linha laranja ) e Nordeste ( linha verde ) apresentam uma situação intermediária. 5.2 CONCEITOS POR CATEGORIA ADMINISTRATIVA, POR MODALIDADE DE ENSINO E GRANDE REGIÃO Na Tabela 5.2, é apresentada a distribuição dos cursos de Zootecnia participantes do Enade/2019 por Categoria Administrativa, de acordo com os conceitos por eles alcançados, segundo as Grandes Regiões brasileiras.
Dos 97 cursos participantes, 83 ( 85,6 \% ) eram ministrados em Instituições Públicas, e 14 ( 14,4 \% ), em Instituições Privadas.
De acordo com as informações da Tabela 5.2, em termos nacionais, entre Instituições Públicas e Privadas, observa -se que, dos quatro cursos avaliados com conceito 5, todos eram oferecidos em IES Públicas.
Dos 83 cursos participantes de IES Públicas, o conceito 3 foi o valor modal, atribuído a 29 cursos ( 34,9 \% ).
Dos demais cursos de Instituições Públicas participantes, oito cursos foram avaliados com conceito 1, 17 cursos ( 20,5 \% ) foram avaliados com conceito 2, 24 cursos ( 28,9 \% ) foram avaliados com conceito 4 e, como já comentado, quatro cursos ( 4,8 \% ), com conceito 5.
Nenhum curso ficou Sem Conceito ( SC ).
Na rede privada, o conceito modal foi 2, atribuído a cinco cursos, o correspondente a 35,7 \% dos 14 cursos da categoria.
Dos demais cursos participantes, dois ( 14,3 \% ) receberam conceito 1, e três ( 21,4 \% ), conceito 3.
O conceito 4 foi atribuído a quatro cursos ( 28,6 \% ).
Nesta Categoria Administrativa, nenhum curso ficou Sem Conceito ( SC ) ou recebeu conceito 5.
Dos 97 cursos participantes da modalidade presencial, o conceito 3 foi o valor modal, atribuído a 32 cursos ( 33,0 \% ).
Dos demais cursos nesta modalidade de ensino, dez cursos ( 10,3 \% ) receberam conceito 1, 22 cursos ( 22,7 \% ) receberam conceito 2, 28 cursos ( 28,9 \% ), conceito 4, e quatro cursos ( 4,1 \% ), conceito 5.
Nesta modalidade de ensino um curso ficou Sem Conceito ( SC ). com conceito 2, 17 cursos ( 20,5 \% ) foram avaliados com conceito 4 e, como já comentado, quatro cursos ( 4,8 \% ), com conceito 5.
Nenhum curso ficou Sem Conceito ( SC ).
Na rede privada, o conceito modal foi 2, atribuído a cinco cursos, o correspondente a 35,7 \% dos 14 cursos da categoria.
Dos demais cursos participantes, dois ( 14,3 \% ) receberam conceito 1, e três ( 21,4 \% ), conceito 3.
O conceito 4 foi atribuído a quatro cursos ( 28,6 \% ).
Nesta Categoria Administrativa, nenhum curso ficou Sem Conceito ( SC ) ou recebeu conceito 5.
Dos 97 cursos participantes da modalidade presencial, o conceito 3 foi o valor modal, atribuído a 32 cursos ( 33,0 \% ).
Dos demais cursos nesta modalidade de ensino, dez cursos ( 10,3 \% ) receberam conceito 1, 22 cursos ( 22,7 \% ) receberam conceito 2, 28 cursos ( 28,9 \% ), conceito 4, e quatro cursos ( 4,1 \% ), conceito 5.
Nesta modalidade de ensino um curso ficou Sem Conceito ( SC ).
Na análise por região, observa -se que, na região Norte, as Instituições Privadas participaram com três cursos ( 23,1 \% do total regional ), dos quais nenhum ficou Sem Conceito ( SC ), e nenhum recebeu os conceitos 2, 4 e 5.
O conceito modal para as Instituições Privadas na região foi 1, atribuído a dois cursos, correspondendo a 66,7 \% dessa categoria na região.
Um curso foi avaliado com conceito 3.
As Instituições Públicas participaram com dez cursos na região Norte ( 76,9 \% do total regional ), aos quais foram atribuídos os conceitos 1, 2, 3 e 4.
Nenhum curso ficou Sem Conceito ( SC ) ou com conceito 5 nessa combinação de categoria e região.
Na região Nordeste, a rede pública concentrou 25 dos 26 cursos participantes, o equivalente a 96,2 \% do total da região.
Dentre os cursos oferecidos por IES Públicas no Nordeste, os conceitos 2 e 3 atribuíram -se a sete cursos cada, sendo os conceitos modais.
Cinco cursos obtiveram conceito 1, quatro cursos receberam conceito 4 e um curso ficou com conceito 5.
Nessa combinação de categoria e região, um curso ficou Sem Conceito ( SC ).
As Instituições Privadas da região Nordeste participaram com um curso ( 3,8 \% do total da região ), que recebeu conceito 2.
Nenhum curso ficou Sem Conceito ( SC ) ou com os demais conceitos nessa região.
Na região Sudeste, a proporção de cursos da rede pública, 79,2 \%, correspondendo a 19 dos 24 cursos participantes.
Nessa categoria e região, o conceito modal foi 3, atribuído a nove cursos.
Os demais foram avaliados com conceito 2 ( um curso ), conceito 4 ( sete cursos ) e conceito 5 ( dois cursos ).
Dos cinco cursos oferecidos por Instituições Privadas na região Sudeste, o conceito modal foi 4, atribuído a três cursos.
Nessa região, nenhum curso de IES Privadas ficou Sem Conceito ( SC ) e nenhum recebeu conceitos 1 e 5.
As Instituições Públicas concentraram 15 dos 16 cursos participantes da região Sul, 93,8 \% do total regional.
Desses, oito cursos obtiveram conceito 4, o conceito modal.
Nesta combinação de Categoria Administrativa e Grande Região, dois cursos receberam conceito 2 e outros cinco, conceito 3.
Nenhum curso recebeu conceito 1 ou 5, e nenhum curso ficou Sem Conceito ( SC ).
As Instituições Privadas na região Sul participaram com um curso ( 6,2 \% dos cursos da região ), ao qual foi atribuído o conceito 4.
Nenhum curso da região Sul oferecido por IES Privadas ficou Sem Conceito ( SC ) ou recebeu os demais conceitos.
Na região Centro-Oeste, 14 dos 18 cursos participantes eram de Instituições Públicas ( 77,8 \% em termos regionais ).
Desses, os conceitos 2 e 3 foram atribuídos a cinco cursos cada, sendo os conceitos modais.
Dos demais cursos, um recebeu conceito 1, outros dois cursos, conceito 4 e um curso, conceito 5.
Nenhum curso ficou Sem Conceito ( SC ).
Quanto aos quatro cursos oferecidos por Instituições Privadas na região Centro-Oeste ( 22,2 \% do total regional ), três cursos apresentaram conceito 2, conceito modal, e um recebeu conceito 3.
Nenhum curso ficou Sem Conceito ( SC ) ou recebeu os demais conceitos nessa combinação de categoria e região.
No Gráfico 5.2, apresenta -se a distribuição do Conceito Enade, segundo a Categoria Administrativa da IES.
Os cursos em IES Públicas ( linha azul ) apresentam uma poligonal mais à direita do que os em IES Privadas ( linha verde ) e, consequentemente, uma distribuição de Conceitos Enade com valores maiores. 5.3 CONCEITOS POR ORGANIZAÇÃO ACADÊMICA E GRANDE REGIÃO Na Tabela 5.3, encontra -se a distribuição dos conceitos atribuídos aos cursos participantes do Enade/2019, na Área de Zootecnia, por Organização Acadêmica, segundo as Grandes Regiões brasileiras.
Dos 97 cursos de Zootecnia participantes, 72 eram oferecidos em Universidades, quatro em Centros Universitários, seis em Faculdades e 15 em CEFET/IF.
Esta distribuição corresponde a, respectivamente, 74,2 \%, 4,1 \%, 6,2 \% e 15,5 \% dos cursos.
De acordo com os dados apresentados, todos os quatro cursos avaliados com conceito 5 eram vinculados a Universidades.
Os cursos vinculados a esse tipo de Organização Acadêmica tiveram o conceito 4 como conceito modal, atribuído a 26 cursos ( 36,1 \% ).
Os demais cursos vinculados a Universidades receberam os conceitos 1 ( cinco cursos ), 2 ( 11 cursos ), 3 ( 25 cursos ) e conceito 5 ( quatro cursos, como já mencionado ).
Um curso ficou Sem Conceito ( SC ).
Dos cursos em Centros Universitários, o conceito modal foi 2, atribuído a dois cursos ( 50,0 \% ).
Nenhum curso vinculado a esse tipo de Organização Acadêmica ficou Sem Conceito ( SC ), e os outros cursos receberam os conceitos 1 ( um curso ) e 4 ( um curso ).
Dos seis cursos mantidos por Faculdades os conceitos 2 e 3 tiveram dois cursos cada, sendo os conceitos modais.
Os conceitos 1 e 4 foram atribuídos a um curso cada.
Dos cursos oferecidos em Faculdades, nenhum ficou Sem Conceito ( SC ) ou recebeu conceito 5.
Dos 15 cursos mantidos por CEFET/IF, sete ( 46,7 \% ) receberam conceito 2, o conceito modal.
O segundo conceito mais frequentemente atribuído aos cursos oferecidos por essa categoria de IES foi o conceito 3, recebido por cinco cursos.
Na sequência, três cursos receberam conceito 1.
Nenhum curso ficou Sem Conceito ( SC ), e nenhum curso recebeu conceito 4 e 5.
Considerando -se, separadamente, as regiões brasileiras, verifica -se que, na região Norte, as Universidades participaram com nove dos 13 cursos avaliados.
Atribuíram -se os conceitos 3 e 4 a três cursos cada, estes os conceitos modais.
Um curso ficou com conceito 1 e dois cursos ficaram com conceito 2.
Nenhum curso ficou Sem Conceito ( SC ) ou recebeu conceito 5.
Os Centros Universitários da região Norte foram representados por um curso, que recebeu conceito 1.
Nenhum curso ficou Sem Conceito ( SC ) ou recebeu os demais conceitos nesta combinação de categoria e região.
As Faculdades participaram com dois cursos na região Norte, os quais receberam os conceitos 1 ( um curso ) e 3 ( um curso ).
Nenhum curso ficou Sem Conceito ( SC ), e nenhum curso oferecido em Faculdades da região obteve conceito 2, 4 ou 5.
Já os CEFET/IF participaram com um curso na região Norte, ao qual foi atribuído o conceito 1.
Na região Nordeste, as Universidades participaram com 21 dos 26 cursos da Área de Zootecnia oferecidos.
O conceito modal 3 foi atribuído sete cursos.
Os outros cursos receberam os conceitos 1 ( três cursos ), 2 ( cinco cursos ), 4 ( quatro cursos ) e 5 ( um curso ).
Um curso oferecido por Universidades no Nordeste ficou Sem Conceito ( SC ).
As Faculdades foram representadas por um curso na região Nordeste, recebendo conceito 2, o conceito modal.
Nenhum curso ficou Sem Conceito ( SC ) ou recebeu os demais conceitos nessa combinação de categoria e região.
Os CEFET/IF participaram com quatro cursos na região Nordeste, avaliados com conceitos 1 ( dois cursos ) e 2 ( também dois cursos ).
Nenhum curso ficou Sem Conceito ( SC ), ou recebeu os conceitos 3, 4 e 5.
Na região Nordeste não houve registro de cursos em Centros Universitários.
Na região Sudeste, as Universidades concentraram 16 dos 24 cursos de Zootecnia da região.
Dos cursos oferecidos em Universidades, nessa região, o conceito modal foi 4, atribuído a oito cursos.
Os demais cursos receberam os conceitos 2 ( um curso ), 3 ( cinco cursos ) e 5 ( dois cursos ).
Nenhum curso ficou Sem Conceito ( SC ) ou recebeu conceito 1.
Os Centros Universitários participaram com dois cursos na região Sudeste, os quais receberam os conceitos 2 e 4.
Nenhum curso ficou Sem Conceito ( SC ) ou recebeu os demais conceitos nessa combinação de categoria e região.
Já as Faculdades foram representadas por dois cursos na região Sudeste, os quais foram atribuídos os conceitos 3 e 4.
Nenhum curso ficou Sem Conceito ( SC ) ou recebeu os demais conceitos nessa combinação de categoria e região.
Os CEFET/IF participaram com quatro cursos, que receberam o conceito 3.
Nenhum curso ficou Sem Conceito ( SC ) ou recebeu os demais conceitos nessa combinação de categoria e região.
Dos 16 cursos da região Sul, 14 eram mantidos por Universidades, dos quais o conceito modal foi 4, atribuído a nove cursos.
Os demais cursos receberam os conceitos 2 ( um curso ) e 3 ( quatro cursos ).
Nesse tipo de organização, nenhum curso da região Sul ficou Sem Conceito ( SC ) ou recebeu conceitos 1 e 5.
Já os CEFET/IF participaram com dois cursos na região Sul, aos quais foram atribuídos os conceitos 2 e 3.
Na região Sul não houve registro de cursos em Centros Universitários e Faculdades.
Na região Centro-Oeste, 12 dos 18 cursos oferecidos eram mantidos por Universidades.
Dos cursos oferecidos em Universidades, nessa região, o conceito modal foi 3, atribuído a seis cursos.
Os demais cursos receberam os conceitos 1 ( um curso ), 2 ( dois cursos ), 4 ( dois cursos ) e 5 ( um curso ).
Nenhum curso ficou Sem Conceito ( SC ).
Os Centros Universitários da região Centro-Oeste contaram com um curso, avaliado com conceito 2.
Um curso foi oferecido em Faculdades na região Centro-Oeste, avaliado com o conceito 2.
Já os CEFET/IF participaram com quatro cursos na região Centro-Oeste, aos quais foram atribuídos o conceito 2.
No Gráfico 5.3, apresenta -se a distribuição do Conceito Enade dos cursos de Zootecnia, segundo a Organização Acadêmica da IES.
Os cursos em Universidades ( linha azul ) apresentam uma poligonal mais à direita do que as demais, denotando uma melhor distribuição de Conceitos Enade.
No outro extremo, com distribuição muito semelhante entre si, encontram -se os cursos mantidos por CEFET/IF ( linha laranja ), Centros Universitários ( linha verde ) e Faculdades ( linha vermelha ), apresentando as respectivas poligonais mais à esquerda e, consequentemente, distribuição de Conceitos Enade com menores valores.

% matched lemmas: 
\end{textsample}
